\usepackage{etex}
\usepackage[italian]{babel}
\usepackage[utf8]{inputenc}
\usepackage[T1]{fontenc}

% Use Moloch theme (evolution of Metropolis)
\usetheme[numbering=fraction,progressbar=frametitle,block=transparent,titleformat=smallcaps]{moloch}
\useoutertheme{moloch}
\useinnertheme{moloch}
\usefonttheme{moloch}
\usecolortheme{moloch}

% Set theme colors
\usepackage[x11names]{xcolor} % See https://latex-tutorial.com/color-latex/ for all loaded colors

\definecolor{venis-darker-blue}{RGB}{2,48,89} % Darker blue extracted with color.adobe.com
\definecolor{venis-lighter-blue}{RGB}{2,72,115} % Lighter blue extracted with color.adobe.com
\definecolor{venis-faithful-blue}{RGB}{0,60,105} % Faithful blue hue extracted from Venis logo
\definecolor{venis-darker-yellow}{RGB}{242,135,5} % Darker yellow hue extracted with color.adobe.com
\definecolor{venis-faithful-yellow}{RGB}{255,161,0} % Faithful yellow hue extracted from Venis logo
\definecolor{monokai-bg}{RGB}{34,35,28} % Faithful Monokai theme for minted package (background)
\definecolor{solarized-light-bg}{RGB}{251,242,200} % Faithful Solarized theme for minted package (background)

\setbeamercolor{normal text}{fg=venis-darker-blue}
\setbeamercolor{alerted text}{fg=venis-darker-yellow}
\setbeamercolor{example text}{fg=venis-faithful-blue}
\setbeamercolor{progress bar}{fg=venis-faithful-yellow}
\setbeamercolor{progress bar in head/foot}{fg=venis-faithful-yellow}
\setbeamercolor{progress bar in section page}{fg=venis-faithful-yellow}
\setbeamercolor{title separator}{fg=venis-darker-yellow}
\setbeamercolor{block title}{bg=venis-lighter-blue,fg=white}
\setbeamercolor{block body}{bg=solarized-light-bg,fg=normal text.fg}


% No navigation symbols on slides
\setbeamertemplate{navigation symbols}{}
% Set the footer
\setbeamertemplate{frame footer}{\alert{\faicon{users}} Simone Scannapieco \qquad\qquad\qquad\qquad\qquad\qquad \alert{\faicon{book}} Spring AI - Corso base \qquad\qquad\qquad\qquad\qquad\qquad \alert{\faicon{institution}} Venis S.p.A, Venezia, IT}
% Set general frame configuration
\setbeamertemplate{page number in head/foot}[appendixframenumber]
\setbeamertemplate{section in toc}[sections numbered]
% Set block configuration
\setbeamertemplate{blocks}[rounded][shadow=true]
% Set special block for code
\newenvironment<>{shellcodeblock}[1]{%
  \setbeamercolor{block body}{bg=solarized-light-bg}%
  \begin{block}{#1}}{\end{block}}

\usepackage[cache=false]{minted}
\usemintedstyle[shell-session]{monokai}
\setminted[shell-session]{
	escapeinside=\@\~\$,
	bgcolor=solarized-light-bg,
	bgcolorpadding=0.5em,
	autogobble,
	fontfamily=tt,
	fontsize=\scriptsize,
	stripall,
	stripnl
}

% Set bibliography properties
\setbeamerfont{bibliography item}{size=\scriptsize}
\setbeamerfont{bibliography entry author}{size=\scriptsize}
\setbeamerfont{bibliography entry title}{size=\scriptsize}
\setbeamerfont{bibliography entry location}{size=\scriptsize}
\setbeamerfont{bibliography entry note}{size=\scriptsize}

\setbeamerfont{block body}{size=\scriptsize}

% Set the color of enumerate labels as the same color of alerted foreground for the current theme
%\setbeamercolor{section number projected}{bg=alerted text.fg}

% Add logo on frame title
\addtobeamertemplate{frametitle}{}{%
	\begin{tikzpicture}[remember picture,overlay]
		\node[anchor=north east,xshift=-5pt,yshift=-7pt] at (current page.north east) 		{\includesvg[height=15pt,keepaspectratio]{dgroove-white.svg}\hspace*{5pt}\includegraphics[height=15pt,keepaspectratio]{venis-logo-round.png}};
	\end{tikzpicture}
  	\ifx\insertframesubtitle\@empty\else%
  		\usebeamerfont{framesubtitle}%
  		\usebeamercolor[fg]{framesubtitle}%
  		\insertframesubtitle%
  	\fi%
}

\usepackage[sfdefault]{roboto}

\usepackage{appendixnumberbeamer}

\usepackage{latexsym}
\usepackage{amssymb,amsfonts}
\usepackage{amsmath}
\usepackage{dsfont}
\usepackage{pifont}

\usepackage{fontspec}
\usepackage{fontawesome}

\usepackage{xspace}
\usepackage[caption=true,font=scriptsize]{subfig}

\usepackage[algo2e,lined,linesnumbered,titlenumbered,longend,italiano]{algorithm2e}

\usepackage{tabularx}
\usepackage{booktabs}
\usepackage{colortbl}
\usepackage{epigraph}

% Multimedia configuration
\usepackage{graphicx}
\graphicspath{{./img/}{../../img/base/}}
% Load the adjustbox package with the export option to make functions available to the \includegraphics command 
\usepackage[export]{adjustbox}
\DeclareGraphicsExtensions{.pdf,.jpeg,.png,.eps}
\usepackage{svg}
\setsvg{inkscapeexe=inkscape, inkscapeopt=-z -D}

% To use fontawesome icons as item label in itemize environments
\usepackage{enumitem}
\usepackage{setspace}
\SetLabelAlign{center}{\strut\smash{\parbox[t]\labelwidth{\centering#1}}}
\setenumerate[1]{%
  label=\protect\usebeamerfont{enumerate item}%
        \protect\usebeamercolor[fg]{enumerate item}%
        \insertenumlabel.
}
\setitemize{label=\usebeamerfont*{itemize item}%
  \usebeamercolor[fg]{itemize item}
  \usebeamertemplate{itemize item}
}

\newcommand*\circled[3]{\tikz[baseline=(char.base)]{						   
            \node[shape=circle,fill=#1,inner sep=1pt] (char)				   
{\color{#2}{\textbf{#3}}};}}

\newcommand*\normalcircled[1]{\circled{venis-darker-blue}{white}{#1}}
\newcommand*\alertedcircled[1]{\circled{venis-darker-yellow}{white}{#1}}

% --------------------------------------------------------------------------------------------------------------------------
% To import media (video and audio):
% (1) STRATEGY 1:
% From: https://tex.stackexchange.com/questions/516029/media9-is-becoming-obsolete-dec-2020-any-alternatives-for-embedding-video-audio/516102#516102

\usepackage{xkeyval}
\usepackage{media9}
\usepackage{zref-abspage}

\makeatletter
\define@boolkey{simplemedia}{autoplay}{\def\simplemediaautoplay{#1}}
\define@boolkey{simplemedia}{showGUI}{\def\simplemediashowGUI{#1}}
\presetkeys{simplemedia}{autoplay=true,showGUI=true}{}

% all page-references relate to: Adobe PDF Reference, sixth edition, November 2006, can be found e.g.:
%https://ghostscript.com/~robin/pdf_reference17.pdf
%%%%%%%%%%%%%%%%%%%%%%%%%%%%%%%%%%%%%%%%%%%%%%%%%%%%%%%%%%%%%%%%%%%%%%%%%%%%%%
% \simplemedia[<options>]{<poster or text>}{<media file>}{MIME type}
%%%%%%%%%%%%%%%%%%%%%%%%%%%%%%%%%%%%%%%%%%%%%%%%%%%%%%%%%%%%%%%%%%%%%%%%%%%%%%
\ExplSyntaxOn
%\cs_new:Npn\simplemedia[3][]{
\newcommand\simplemedia[4][]{%
  \setkeys{simplemedia}{#1}%
  \leavevmode
  \pbs_pdfobj:nnn{}{fstream}{{}{#3}}
  \pbs_pdfobj:nnn{}{dict}{
    /Type/Filespec/F~(#3)/UF~(#3)
    /EF~<</F~\pbs_pdflastobj:>>
  }
  %
  \pbs_pdfobj:nnn{}{dict}{
      /CT~(#4)% content type, page 764, table 9.9
      /D~\pbs_pdflastobj:% full file spcification, page 764, table 9.9
      /N~(Media~clip~from~#3)% name of media clip, page 764, table 9.8
      /P~<</TF (TEMPACCESS)>>% media permissions, page 764, table 9.9
      /S~/MCD% subtype media-clip-data, page 764, table 9.8
  }
  %
  \pbs_pdfobj:nnn{}{dict}{
      /C~\pbs_pdflastobj:% media clip dictionary, page 762, table 9.6
      /N~(RenditionFrom#3)% name of rendition, page 759, table 9.1
      /P~<<% media play parameters dictionary, page 762, table 9.6
%%%%%%%%%%%%%%%%%%%%%%%%%%%%%%%%%%%%%%%%%%%%%%%%%%%%%%%%%%%%%%%%%%%%%%%%%%%%%%%%%%%%%%%%%%%%%%%%%%%%%%%%%%%%%%%%%%%%%%%%%%
%       Here's the point to switch on/off player GUI
%%%%%%%%%%%%%%%%%%%%%%%%%%%%%%%%%%%%%%%%%%%%%%%%%%%%%%%%%%%%%%%%%%%%%%%%%%%%%%%%%%%%%%%%%%%%%%%%%%%%%%%%%%%%%%%%%%%%%%%%%%
        /BE~<</C~\simplemediashowGUI~%C: display a player-specific controller user interface, page 769, table 9.15
        /A~true>>~%A: automatically play when activated, page 770, table 9.15
        % select player (not used here)
%        /PL~<</A~[% arrray of media player info objects page 777
%%        <</PID~<</U~(vnd.adobe.swname:ADBE_MCI)~>>~>>~% PID: software identifier object, page 779
%        <</PID~<</U~(vnd.adobe.swname:MSFT_WindowsMediaPlayer)~/OS[(win16)~(win32)~(win9x)~(winnt)~(wince)]>>~>>~%/OS operating system, page 780
%        ]~>>
      ~>>
      /S~/MR% subtype of rendition, page 759, table 9.1
  }
  %
  \tl_set:Nx\mediarendition{\pbs_pdflastobj:}
  %  
  \pbs_pdfobj:nnn{}{dict}{}% obj-number created for later use
  \tl_set:Nx\renditionactionplay{\pbs_pdflastobj:}%
  %
  \pbs_pdfobj:nnn{}{dict}{}% obj-number created for later use
  \tl_set:Nx\renditionactionstop{\pbs_pdflastobj:}%
  %
  \hbox_set:Nn\l_tmpa_box{#2}% size of annotation
  \tl_set:Nx\l_box_wd_tl{\dim_use:N\box_wd:N\l_tmpa_box}
  \tl_set:Nx\l_box_ht_tl{\dim_use:N\box_ht:N\l_tmpa_box}
  \tl_set:Nx\l_box_dp_tl{\dim_use:N\box_dp:N\l_tmpa_box}
  \pbs_pdfxform:nnnnn{1}{1}{}{}{\l_tmpa_box}
  %
  \tl_set:Nx\l_simplem_current_page_tl{\zref@extract{simplem@\int_use:N\g_simplem_id_int}{abspage}}
  \pbs_pdfannot:nnnn{\l_box_wd_tl}{\l_box_ht_tl}{\l_box_dp_tl}{
      /T~(AnnotationFrom#3)~% title of screen annotation, page 640, table 8.38 
      /Subtype~/Screen~% Subtype Sreen, page 616, table 8.20 and page 640, table 8.38
      /BS~<</W~1/S/S>>% Borderstyle
      /C~[0.039216~0.039216~0.039216]% color used background, title bar, border, page 607, table 8.15
      /Contents~(EmbeddedAudiofile#3)% Text to be displayed, page 606, table 8.15
      /NM~(audioscreenannot:#3)% annotation name, page 606, table 8.15
      /AP~<</N~\pbs_pdflastxform:>>% appearance dictionary, page 606, table 8.15
      /F~4~%Annotation flag, page 606, table 8.15 and page 608, table 8.16
      /A~\renditionactionplay~% Action to be performed, page 640 
%%%%%%%%%%%%%%%%%%%%%%%%%%%%%%%%%%%%%%%%%%%%%%%%%%%%%%%%%%%%%%%%%%%%%%%%%%%%%%%%%%%%%%%%%%%%%%%%%%%%%%%%%%%%%%%%%%%%%%%%%%
%       Here's the point to switch on/off autoplay
%%%%%%%%%%%%%%%%%%%%%%%%%%%%%%%%%%%%%%%%%%%%%%%%%%%%%%%%%%%%%%%%%%%%%%%%%%%%%%%%%%%%%%%%%%%%%%%%%%%%%%%%%%%%%%%%%%%%%%%%%%
      % autoplay is realised by "additional action" (page 640), 
      % PageOpen (PO) -> play "Play Rendition"
      % PageClose (PC) -> play "Stop Rendition" (see page 650)
      /AA~<<\ifKV@simplemedia@autoplay~/PO~\renditionactionplay~\fi~/PC~\renditionactionstop~>>~%
%%%%%%%%%%%%%%%%%%%%%%%%%%%%%%%%%%%%%%%%%%%%%%%%%%%%%%%%%%%%%%%%%%%%%%%%%%%%%%%%%%%%%%%%%%%%%%%%%%%%%%%%%%%%%%%%%%%%%%%%%%
%       see Adobe pdfmark-reference, I THINK, THIS WORKS ONLY WITH tex->dvips->ps2pdf
%%%%%%%%%%%%%%%%%%%%%%%%%%%%%%%%%%%%%%%%%%%%%%%%%%%%%%%%%%%%%%%%%%%%%%%%%%%%%%%%%%%%%%%%%%%%%%%%%%%%%%%%%%%%%%%%%%%%%%%%%%
      \int_compare:nT{\l_simplem_current_page_tl>0}{
        /P~\exp_args:Ne\pdf_pageobject_ref:n{\l_simplem_current_page_tl}
      }
      /CA~1%constant opacity value, page 618
  }
  % Rendition Action: "Play Rendition"
  \pbs_pdfobj:nnn{\renditionactionplay}{dict}{%
      /R~\mediarendition%
      /S~/Rendition
      /OP~4%operation to perform (page 669): 4 = play rendition specified by R
      /AN~\pbs_pdflastann:%
  }
  % Rendition Action: "Stop Rendition"
  \pbs_pdfobj:nnn{\renditionactionstop}{dict}{%
      /R~\mediarendition%
      /S~/Rendition
      /OP~1%operation to perform (page 669): 1 = stop rendition specified by R
      /AN~\pbs_pdflastann:%
  }
  %
  \phantom{#2}
  \zref@labelbyprops{simplem@\int_use:N\g_simplem_id_int}{abspage}
  \int_gincr:N\g_simplem_id_int
}%
\int_new:N\g_simplem_id_int
\ExplSyntaxOff
\makeatother

% (2) STRATEGY 2:
% From https://www.youtube.com/watch?v=BspamyoyjUk and https://gist.github.com/FedericoTartarini/7af4eb6fc13b1cb9cc68b7e8ea823d50
% This code was copied from https://tex.stackexchange.com/a/516102

\usepackage[bigfiles]{pdfbase}
\ExplSyntaxOn
\NewDocumentCommand\embedvideo{smm}{
  \group_begin:
  \leavevmode
  \tl_if_exist:cTF{file_\file_mdfive_hash:n{#3}}{
    \tl_set_eq:Nc\video{file_\file_mdfive_hash:n{#3}}
  }{
    \IfFileExists{#3}{}{\GenericError{}{File~`#3'~not~found}{}{}}
    \pbs_pdfobj:nnn{}{fstream}{{}{#3}}
    \pbs_pdfobj:nnn{}{dict}{
      /Type/Filespec/F~(#3)/UF~(#3)
      /EF~<</F~\pbs_pdflastobj:>>
    }
    \tl_set:Nx\video{\pbs_pdflastobj:}
    \tl_gset_eq:cN{file_\file_mdfive_hash:n{#3}}\video
  }
  %
  \pbs_pdfobj:nnn{}{dict}{
    /Type/RichMediaInstance/Subtype/Video
    /Asset~\video
    /Params~<</FlashVars (
      source=#3&
      skin=SkinOverAllNoFullNoCaption.swf&
      skinAutoHide=true&
      skinBackgroundColor=0x5F5F5F&
      skinBackgroundAlpha=0
    )>>
  }
  %
  \pbs_pdfobj:nnn{}{dict}{
    /Type/RichMediaConfiguration/Subtype/Video
    /Instances~[\pbs_pdflastobj:]
  }
  %
  \pbs_pdfobj:nnn{}{dict}{
    /Type/RichMediaContent
    /Assets~<<
      /Names~[(#3)~\video]
    >>
    /Configurations~[\pbs_pdflastobj:]
  }
  \tl_set:Nx\rmcontent{\pbs_pdflastobj:}
  %
  \pbs_pdfobj:nnn{}{dict}{
    /Activation~<<
      /Condition/\IfBooleanTF{#1}{PV}{XA}
      /Presentation~<</Style/Embedded>>
    >>
    /Deactivation~<</Condition/PI>>
  }
  %
  \hbox_set:Nn\l_tmpa_box{#2}
  \tl_set:Nx\l_box_wd_tl{\dim_use:N\box_wd:N\l_tmpa_box}
  \tl_set:Nx\l_box_ht_tl{\dim_use:N\box_ht:N\l_tmpa_box}
  \tl_set:Nx\l_box_dp_tl{\dim_use:N\box_dp:N\l_tmpa_box}
  \pbs_pdfxform:nnnnn{1}{1}{}{}{\l_tmpa_box}
  %
  \pbs_pdfannot:nnnn{\l_box_wd_tl}{\l_box_ht_tl}{\l_box_dp_tl}{
    /Subtype/RichMedia
    /BS~<</W~0/S/S>>
    /Contents~(embedded~video~file:#3)
    /NM~(rma:#3)
    /AP~<</N~\pbs_pdflastxform:>>
    /RichMediaSettings~\pbs_pdflastobj:
    /RichMediaContent~\rmcontent
  }
  \phantom{#2}
  \group_end:
}
\ExplSyntaxOff
% --------------------------------------------------------------------------------------------------------------------------

\mode<beamer|second|trans|article>{
  \newcommand\presentvideo[2]{\embedvideo{\includegraphics[width=\textwidth]{#1}}{#2}}
}

\mode<handout>{
  \newcommand\presentvideo[2]{\includegraphics[width=\textwidth]{#1}}
}

\makeatletter
\def\beamer@framenotesbegin{% at beginning of slide
\usebeamercolor[fg]{normal text}
\gdef\beamer@noteitems{}%
\gdef\beamer@notes{}%
}
\makeatother

\usepackage{hyperref}

\hypersetup{
    colorlinks=true,
    linkcolor=venis-faithful-blue,
    filecolor=venis-faithful-yellow,      
    urlcolor=venis-faithful-blue,
    pdftitle={Overleaf Example},
    pdfpagemode=FullScreen,
    }